% GENERATED FILE. Edit canonical lessons, then run npm run generate:books.
% canonical-source-hash: fnv1a64:ec71650779225e99
% canonical-lessons: RU-C120-lion, RU-C120-monkey, RU-C120-elephant, RU-C120-bee, RU-C120-ant

\chapter{Lion, Monkey, Elephant, Bee, Ant}
\label{ch:ru-lion-monkey-elephant-bee-ant}
\hlchaptermodality{120}

\begin{chapteropening}
\textbf{By the end of this chapter:} \emph{I can say a lion, a monkey, an elephant, a bee and an ant.}

Complete the last lesson of chapter 120: I can say a lion, a monkey, an elephant, a bee and an ant.
\end{chapteropening}

\section[\texorpdfstring{\ru{лев} (lev)}{лев (lev)}]{\ru{лев} (lev) — a lion}

\label{lesson:RU-C120-lion}



\begin{quote}
Before the new one: say the Russian for lips, then the Russian for the chest.
\end{quote}

\subsection*{You'll want to know: \ru{лев}}
\textbf{\ru{лев}} (\emph{lev}) — \textquotedblleft{}a lion\textquotedblright{}.

\begin{cousinweb}
\textbf{\ru{лев}} $\leftarrow$ Old High German \emph{lëwo}, from Latin \emph{leō}, Greek \emph{léōn} (Vasmer). Vasmer thinks it came through books, since a spoken loan would have given \emph{\ru{лёв}}, and that the word's ultimate source is probably Semitic.
\end{cousinweb}

\begin{grammarlens}[title={What you've built}]
One more word for this chapter.
\end{grammarlens}

\subsection*{Your turn}
\begin{itemize}
\raggedright
  \item \emph{Say it:} \emph{lev}
  \item \emph{Say it:} \emph{lev}, once more
  \item \emph{Say it:} say \emph{lev}
  \item \emph{Recall:} say the Russian for lips, then the Russian for the chest, then say \emph{lev} again
\end{itemize}

\begin{tcolorbox}[breakable,colback=teal!4,colframe=teal!35!black,title={Before you move on}]
What does \ru{лев} mean? (\textquotedblleft{}A lion\textquotedblright{}.) Say it once more.
\end{tcolorbox}
\section[\texorpdfstring{\ru{обезьяна} (obezyána)}{обезьяна (obezyána)}]{\ru{обезьяна} (obezyána) — a monkey}

\label{lesson:RU-C120-monkey}



\begin{quote}
Before the new one: say the Russian for the chest, then the Russian for a lion.
\end{quote}

\subsection*{You'll want to know: \ru{обезьяна}}
\textbf{\ru{обезьяна}} (\emph{obezyána}) — \textquotedblleft{}a monkey\textquotedblright{}.

\begin{cousinweb}
\textbf{\ru{обезьяна}} — Old Russian \emph{\ru{обезьяна}}, recorded by Afanasy Nikitin; it comes from Turkish and Persian \emph{abuzine}, \textquotedblleft{}monkey\textquotedblright{}, reshaped to look like Russian words with \emph{\ru{о}-} and \emph{\ru{без}-} (Vasmer).
\end{cousinweb}

\begin{grammarlens}[title={What you've built}]
One more word for this chapter.
\end{grammarlens}

\subsection*{Your turn}
\begin{itemize}
\raggedright
  \item \emph{Say it:} \emph{obezyána}
  \item \emph{Say it:} \emph{obezyána}, once more
  \item \emph{Say it:} say \emph{obezyána}
  \item \emph{Recall:} say the Russian for the chest, then the Russian for a lion, then say \emph{obezyána} again
\end{itemize}

\begin{tcolorbox}[breakable,colback=teal!4,colframe=teal!35!black,title={Before you move on}]
What does \ru{обезьяна} mean? (\textquotedblleft{}A monkey\textquotedblright{}.) Say it once more.
\end{tcolorbox}
\section[\texorpdfstring{\ru{слон} (slon)}{слон (slon)}]{\ru{слон} (slon) — an elephant}

\label{lesson:RU-C120-elephant}



\begin{quote}
Before the new one: say the Russian for a lion, then the Russian for a monkey.
\end{quote}

\subsection*{You'll want to know: \ru{слон}}
\textbf{\ru{слон}} (\emph{slon}) — \textquotedblleft{}an elephant\textquotedblright{}.

\begin{cousinweb}
\textbf{\ru{слон}} — origin disputed. It is usually read as a bookish coinage from \emph{\ru{прислониться}}, \textquotedblleft{}to lean\textquotedblright{}, after the old \emph{Physiologus}, which says the elephant sleeps leaning on a tree; Vasmer thinks that link is folk etymology and finds a loan from Turkic \emph{arslan}, \textquotedblleft{}lion\textquotedblright{}, acceptable in its sounds.
\end{cousinweb}

\begin{grammarlens}[title={What you've built}]
One more word for this chapter.
\end{grammarlens}

\subsection*{Your turn}
\begin{itemize}
\raggedright
  \item \emph{Say it:} \emph{slon}
  \item \emph{Say it:} \emph{slon}, once more
  \item \emph{Say it:} say \emph{slon}
  \item \emph{Recall:} say the Russian for a lion, then the Russian for a monkey, then say \emph{slon} again
\end{itemize}

\begin{tcolorbox}[breakable,colback=teal!4,colframe=teal!35!black,title={Before you move on}]
What does \ru{слон} mean? (\textquotedblleft{}An elephant\textquotedblright{}.) Say it once more.
\end{tcolorbox}
\section[\texorpdfstring{\ru{пчела} (pchelá)}{пчела (pchelá)}]{\ru{пчела} (pchelá) — a bee}

\label{lesson:RU-C120-bee}



\begin{quote}
Before the new one: say the Russian for a monkey, then the Russian for an elephant.
\end{quote}

\subsection*{You'll want to know: \ru{пчела}}
\textbf{\ru{пчела}} (\emph{pchelá}) — \textquotedblleft{}a bee\textquotedblright{}.

\begin{cousinweb}
\textbf{\ru{пчела}} $\leftarrow$ Proto-Slavic \emph{*bĭčela} — inherited; \emph{\ru{б}} turned to \emph{\ru{п}} before \emph{\ru{ч}}. Vasmer relates it to Irish \emph{bech}, Lithuanian \emph{bìtė} and Old High German \emph{bîa}, all \textquotedblleft{}bee\textquotedblright{}: English \emph{bee} is a cousin.
\end{cousinweb}

\begin{grammarlens}[title={What you've built}]
One more word for this chapter.
\end{grammarlens}

\subsection*{Your turn}
\begin{itemize}
\raggedright
  \item \emph{Say it:} \emph{pchelá}
  \item \emph{Say it:} \emph{pchelá}, once more
  \item \emph{Say it:} say \emph{pchelá}
  \item \emph{Recall:} say the Russian for a monkey, then the Russian for an elephant, then say \emph{pchelá} again
\end{itemize}

\begin{tcolorbox}[breakable,colback=teal!4,colframe=teal!35!black,title={Before you move on}]
What does \ru{пчела} mean? (\textquotedblleft{}A bee\textquotedblright{}.) Say it once more.
\end{tcolorbox}
\section[\texorpdfstring{\ru{муравей} (muravéy)}{муравей (muravéy)}]{\ru{муравей} (muravéy) — an ant}

\label{lesson:RU-C120-ant}



\begin{quote}
Before the new one: say the Russian for an elephant, then the Russian for a bee.
\end{quote}

\subsection*{You'll want to know: \ru{муравей}}
\textbf{\ru{муравей}} (\emph{muravéy}) — \textquotedblleft{}an ant\textquotedblright{}.

\begin{cousinweb}
\textbf{\ru{муравей}} $\leftarrow$ Proto-Slavic \emph{*morvĭjĭ} — inherited. Vasmer lists Irish \emph{moirb}, Avestan \emph{maoiri-} and Old Norse \emph{maurr}, all \textquotedblleft{}ant\textquotedblright{}; Latin \emph{formīca}, \textquotedblleft{}ant\textquotedblright{}, is related too.
\end{cousinweb}

\begin{grammarlens}[title={What you've built}]
One more word for this chapter.
\end{grammarlens}

\subsection*{Your turn}
\begin{itemize}
\raggedright
  \item \emph{Say it:} \emph{muravéy}
  \item \emph{Say it:} \emph{muravéy}, once more
  \item \emph{Say it:} say \emph{muravéy}
  \item \emph{Recall:} say the Russian for an elephant, then the Russian for a bee, then say \emph{muravéy} again
\end{itemize}

\begin{tcolorbox}[breakable,colback=teal!4,colframe=teal!35!black,title={Before you move on}]
What does \ru{муравей} mean? (\textquotedblleft{}An ant\textquotedblright{}.) Say it once more.
\end{tcolorbox}
